% Focused framework presentation.
\section[Algebraic solution]{Solving dynamics with evaluation and kernels}

\begin{frame}[t]{An affine ODE as a linear operator equation}
\hypertarget{main:ivp}{}
\begin{columns}[T,onlytextwidth]
  \begin{column}{.54\textwidth}
    \heading{Given data}
    An open interval $I\subseteq\mathbb R$, with
    \begin{equation*}
      \begin{aligned}
        A&\in C(I,\operatorname{End}(V)),\\
        B&\in C(I,V).
      \end{aligned}
    \end{equation*}
    For $(s,x)\in I\times V$, solve
    \begin{equation*}
      \begin{cases}
        \dot x_t=A_tx_t+B_t,&t\in I,\\
        x_s=x.
      \end{cases}
    \end{equation*}
    Continuous coefficients give a unique trajectory on all of $I$.
  \end{column}
  \begin{column}{.41\textwidth}
    \heading{Trajectory formulation}
    \flowdiagram{trajectory}
    \begin{equation*}
      \mathsf D_A\gamma=B,\qquad E_s\gamma=x.
    \end{equation*}
    The unknown is a complete trajectory.
  \end{column}
\end{columns}
\begin{equation*}
  \mathcal X=C^1(I,V),\qquad \mathcal Y=C(I,V),\qquad
  (\mathsf M_A\gamma)_t=A_t\gamma_t.
\end{equation*}
\slidecite{Assessment 9.1; Manuscript Example 1.10. \hyperlink{app:existence}{Existence and uniqueness: Appendix A2.}}
\end{frame}

\begin{frame}[t]{Homogeneous evolution by evaluation}
\hypertarget{main:homogeneous}{}
\begin{columns}[T,onlytextwidth]
  \begin{column}{.54\textwidth}
    \heading{The homogeneous kernel}
    \begin{equation*}
      \mathcal H=\ker\mathsf D_A.
    \end{equation*}
    Evaluation identifies a trajectory with its state:
    \begin{equation*}
      E_s:\mathcal H\to V,\qquad E_s(\eta)=\eta_s.
    \end{equation*}
    Existence and uniqueness make $E_s$ a linear isomorphism.
    \begin{equation*}
      S_s=E_s^{-1},\qquad U(t,s)=E_tS_s.
    \end{equation*}
  \end{column}
  \begin{column}{.41\textwidth}
    \heading{Recover, then evaluate}
    \flowdiagram{linear-evolution}
    \medskip
    Write $U_t=U(t,s)$. Then
    \begin{equation*}
      \begin{aligned}
        \dot U_t&=A_tU_t,\\
        U_s&=I_V.
      \end{aligned}
    \end{equation*}
  \end{column}
\end{columns}
\slidecite{Manuscript Example 1.10. \hyperlink{app:evaluation}{Evaluation isomorphism and composition: Appendix A3.}}
\end{frame}

\begin{frame}[t]{The affine solution and its flow}
\hypertarget{main:affine}{}
\begin{columns}[T,onlytextwidth]
  \begin{column}{.55\textwidth}
    \heading{Homogeneous and forced response}
    \begin{equation*}
      \begin{aligned}
        x_t&=U(t,s)x+a_{s,t},\\
        a_{s,t}&=\int_s^t U(t,r)B_r\,\mathrm dr.
      \end{aligned}
    \end{equation*}
    The initial value selects one trajectory from an affine family.
    \heading{Initial-to-terminal state map}
    \begin{equation*}
      \Phi_{s,t}(x)=U(t,s)x+a_{s,t}.
    \end{equation*}
    For zero forcing, the map reduces to $U(t,s)x$.
  \end{column}
  \begin{column}{.40\textwidth}
    \flowdiagram[.92\linewidth]{affine}
    \begin{equation*}
      \mathcal S_B=\gamma_{\mathrm p}+\mathcal H.
    \end{equation*}
  \end{column}
\end{columns}
\slidecite{Manuscript Example 1.10. \hyperlink{app:affine}{Affine family: A4}; \hyperlink{app:variation}{derivation: A5}; \hyperlink{app:verification}{verification: A6}.}
\end{frame}

\begin{frame}[t]{Evaluation also constructs nonlinear flows}
\hypertarget{main:general}{}
\begin{columns}[T,onlytextwidth]
  \begin{column}{.55\textwidth}
    \heading{Allow a nonlinear velocity field}
    \begin{equation*}
      \dot x_t=v_t(x_t),\qquad x_s=x,
      \qquad \Phi_{s,t}(x)=x_t.
    \end{equation*}
    Assume continuous, locally Lipschitz dynamics with a unique trajectory on the same interval $I$ for every initial pair.
    \heading{The algebraic construction remains}
    \begin{equation*}
      \Phi_{s,t}=E_t\circ E_s^{-1}.
    \end{equation*}
    Evaluation is a bijection of sets; linearity of the state map is unnecessary.
  \end{column}
  \begin{column}{.40\textwidth}
    \vspace{8pt}
    \flowdiagram{general-flow}
  \end{column}
\end{columns}
\slidecite{Assessment 9.2; Manuscript Theorem 1.12. \hyperlink{app:general}{General evaluation construction: Appendix A7.}}
\end{frame}

\begin{frame}[t]{Flow laws by cancellation of evaluation maps}
\hypertarget{main:laws}{}
On the set $\mathcal S_v$ of complete solution trajectories, existence and
uniqueness make each evaluation $E_s:\mathcal S_v\to V$ a bijection.
\heading{Composition follows from an inverse identity}
\begin{equation*}
  \begin{aligned}
    \Phi_{r,t}\circ\Phi_{s,r}
      &=(E_t\circ E_r^{-1})\circ(E_r\circ E_s^{-1})\\
      &=E_t\circ E_s^{-1}=\Phi_{s,t}.
  \end{aligned}
\end{equation*}
\begin{columns}[T,onlytextwidth]
  \begin{column}{.48\textwidth}
    \heading{Identity and inverse}
    \begin{equation*}
      \Phi_{s,s}=I_V,\qquad \Phi_{s,t}^{-1}=\Phi_{t,s}.
    \end{equation*}
    The identities hold on a common interval for all initial data.
  \end{column}
  \begin{column}{.47\textwidth}
    \flowdiagram{flow-composition}
  \end{column}
\end{columns}
\medskip
The flow identities use composition and invertibility of evaluation.
Linearity of the state map is unnecessary.
\slidecite{Manuscript \S\,1.2.1; Assessment 9.2. \hyperlink{app:general}{Evaluation construction: Appendix A7.}}
\end{frame}
